% packages
%\usepackage[colorlinks,linkcolor=blue,citecolor=blue]{hyperref}
%\usepackage{amsmath,amsfonts,amssymb}
%\usepackage{mathtools}
%\usepackage{amsthm} % better to load after ams math
%\usepackage{latexsym}

%%\usepackage{baskervald}
%\usepackage[p,theoremfont]{baskervillef}
%\usepackage[baskervillef,smallerops,vvarbb]{newtxmath}
%%\usepackage[cal=boondoxo]{mathalfa}

%\usepackage{newtxtext}
%\usepackage[smallerops,vvarbb]{newtxmath}
%\usepackage[cal=boondoxo]{mathalfa}

%\usepackage{mathabx}
%\usepackage{newtxtext}
%%\usepackage{libertine}
%%\usepackage[libertine]{newtxmath}
%\usepackage[noamssymbols]{newtxmath}
%\usepackage{relsize}
%\usepackage{thmtools,thm-restate}
%\usepackage[amsthm]{ntheorem}

\usepackage{enumitem}
\usepackage{bm}
\usepackage{xspace}
\usepackage{nicefrac}
%\usepackage{keyval}
%\usepackage{ifthen}

% allow using \bm in section titles
\pdfstringdefDisableCommands{%
  \renewcommand{\bm}[1]{#1}%
}

\usepackage[capitalise,nameinlink]{cleveref}
\usepackage{xcolor}
\usepackage{dsfont}


% comments, etc
\newcommand{\ignore}[1]{}
\newcommand{\tk}[1]{\noindent{\textcolor{cyan}{\{{\bf TK:} {\em #1}\}}}}
% for emails
\newcommand{\email}[1]{\texttt{#1}}

%% JMLR theorems adjustments
%\setlength{\topsep}{\medskipamount}
%
%\newtheoremstyle{jmlrthm} % name
%	{\topsep}		% Space above
%	{\topsep}		% Space below
%	{\itshape}	% Body font
%	{}						% Indent amount
%	{\bfseries}	% Theorem head font
%	{}						% Punctuation after theorem head
%	{0.33em}		% Space after theorem head
%	{}  					% Theorem head spec (can be left empty, meaning ?normal?)
%
%\newtheoremstyle{jmlrdef} % name
%	{\topsep}		% Space above
%	{\topsep}		% Space below
%	{}						% Body font
%	{}						% Indent amount
%	{\bfseries}	% Theorem head font
%	{}						% Punctuation after theorem head
%	{.33em}			% Space after theorem head
%	{}  					% Theorem head spec (can be left empty, meaning ?normal?)
%
%\renewcommand\qedsymbol{\BlackBox}
%%\renewenvironment{proof}[1][Proof]{\par\noindent{\bfseries\upshape #1\ }}{\hfill\BlackBox\\[\medskipamount]}
%\renewenvironment{proof}[1][Proof]{\par\noindent{\bfseries\upshape #1\ }}{\hfill\BlackBox\\}

% wrap algorithm
\newcommand{\wrapalgo}[2][0.9\linewidth]
{%
\begin{center}\setlength{\fboxsep}{5pt}\fbox{\begin{minipage}{#1}
#2
\end{minipage}}\end{center}
\vspace{-0.25cm}
}


% smash text (for inline math)
\newenvironment{smashbox}[1][]{%
  \begin{minipage}{1.0\linewidth}\lineskiplimit=-\maxdimen
}{%
  \end{minipage}
}



% theorems
%\theoremstyle{plain}
%\theoremstyle{jmlrthm}
\theoremstyle{definition}
\newtheorem{theorem}{Theorem}				% global numbering
%\newtheorem{theorem}{Theorem}[section]		% section numbering

\newtheorem{lemma}[theorem]{Lemma}
\newtheorem{corollary}[theorem]{Corollary}
\newtheorem{proposition}[theorem]{Proposition}
\newtheorem{claim}[theorem]{Claim}
\newtheorem{fact}[theorem]{Fact}
\newtheorem{observation}[theorem]{Observation}

\newtheorem*{theorem*}{Theorem}
\newtheorem*{lemma*}{Lemma}
\newtheorem*{corollary*}{Corollary}
\newtheorem*{proposition*}{Proposition}
\newtheorem*{claim*}{Claim}
\newtheorem*{fact*}{Fact}
\newtheorem*{observation*}{Observation}

\theoremstyle{definition}
%\theoremstyle{jmlrdef}
\newtheorem{definition}[theorem]{Definition}
\newtheorem{remark}[theorem]{Remark}
\newtheorem{example}[theorem]{Example}
\newtheorem{question}[theorem]{Question}
\newtheorem*{definition*}{Definition}
\newtheorem*{remark*}{Remark}
\newtheorem*{example*}{Example}
\newtheorem*{question*}{Question}


% repeated theorems
 \theoremstyle{plain}
%\theoremstyle{jmlrthm}
\newtheorem*{theoremaux}{\theoremauxref}
\gdef\theoremauxref{1}

\newenvironment{repthm}[2][]{%
  \def\theoremauxref{\cref{#2}}
  \begin{theoremaux}[#1]
}{%
  \end{theoremaux}
}

%\newtheoremstyle{theoremrep}
%	{\topsep}{\topsep}              %%% space between body and thm
%	{\itshape}                      %%% Thm body font
%	{}                              %%% Indent amount (empty = no indent)
%	{\bfseries}                     %%% Thm head font
%	{}                             %%% Punctuation after thm head
%	{0.33em}                             %%% Space after thm head
%	{\cref{#3}}
%
%\theoremstyle{theoremrep}
%\newtheorem{repthm}{dontcare}



% % old refs - use \cref{} instead!
% \newcommand{\secref}[1]{Section~\ref{#1}}
%\newcommand{\secref}[1]{Sec.~\ref{#1}}
%\newcommand{\figref}[1]{Figure~\ref{#1}}
%\renewcommand{\eqref}[1]{Eq.~(\ref{#1})}
%\newcommand{\lemref}[1]{Lemma~\ref{#1}}

% \newcommand{\thmref}[1]{Theorem~\ref{#1}}
% \newcommand{\appref}[1]{Appendix~\ref{#1}}
% \newcommand{\algref}[1]{Algorithm~\ref{#1}}

% style defs
%\renewcommand{\le}{\leqslant}
%\renewcommand{\ge}{\geqslant}
%\renewcommand{\leq}{\le}
%\renewcommand{\geq}{\ge}

% bold serif
\DeclareMathAlphabet{\mathbfsf}{\encodingdefault}{\sfdefault}{bx}{n}

% operators
\DeclareMathOperator*{\err}{error}
\DeclareMathOperator*{\argmin}{arg\!\min}
\DeclareMathOperator*{\argmax}{arg\!\max}
\DeclareMathOperator*{\conv}{conv}
\DeclareMathOperator*{\supp}{supp}
\DeclareMathOperator*{\trace}{Tr}
\DeclareMathOperator*{\diag}{diag}
%\DeclareMathOperator*{\poly}{poly}
\let\Pr\relax
\DeclareMathOperator{\Pr}{\mathbb{P}}

% cases
\newcommand{\mycases}[4]{{
\left\{
\begin{array}{ll}
    {#1} & {\;\text{#2}} \\[1ex]
    {#3} & {\;\text{#4}}
\end{array}
\right. }}

\newcommand{\mythreecases}[6] {{
\left\{
\begin{array}{ll}
    {#1} & {\;\text{#2}} \\[1ex]
    {#3} & {\;\text{#4}} \\[1ex]
    {#5} & {\;\text{#6}}
\end{array}
\right. }}

% macros
\newcommand{\lr}[1]{\mathopen{}\left(#1\right)}
\newcommand{\Lr}[1]{\mathopen{}\big(#1\big)}
\newcommand{\LR}[1]{\mathopen{}\Big(#1\Big)}

\newcommand{\lrbra}[1]{\mathopen{}\left[#1\right]}
\newcommand{\Lrbra}[1]{\mathopen{}\big[#1\big]}
\newcommand{\LRbra}[1]{\mathopen{}\Big[#1\Big]}

\newcommand{\norm}[1]{\|#1\|}
\newcommand{\lrnorm}[1]{\mathopen{}\left\|#1\right\|}
\newcommand{\Lrnorm}[1]{\mathopen{}\big\|#1\big\|}
\newcommand{\LRnorm}[1]{\mathopen{}\Big\|#1\Big\|}

\newcommand{\set}[1]{\{#1\}}
\newcommand{\lrset}[1]{\mathopen{}\left\{#1\right\}}
\newcommand{\Lrset}[1]{\mathopen{}\big\{#1\big\}}
\newcommand{\LRset}[1]{\mathopen{}\Big\{#1\Big\}}

\newcommand{\abs}[1]{|#1|}
\newcommand{\lrabs}[1]{\mathopen{}\left|#1\right|}
\newcommand{\Lrabs}[1]{\mathopen{}\big|#1\big|}
\newcommand{\LRabs}[1]{\mathopen{}\Big|#1\Big|}

\newcommand{\ceil}[1]{\lceil #1 \rceil}
\newcommand{\lrceil}[1]{\mathopen{}\left\lceil #1 \right\rceil}
\newcommand{\Lrceil}[1]{\mathopen{}\big\lceil #1 \big\rceil}
\newcommand{\LRceil}[1]{\mathopen{}\Big\lceil #1 \Big\rceil}

\newcommand{\floor}[1]{\lfloor #1 \rfloor}
\newcommand{\lrfloor}[1]{\mathopen{}\left\lfloor #1 \right\rfloor}
\newcommand{\Lrfloor}[1]{\mathopen{}\big\lfloor #1 \big\rfloor}
\newcommand{\LRfloor}[1]{\mathopen{}\Big\lfloor #1 \Big\rfloor}

\newcommand{\ip}[1]{\langle #1 \rangle}
\newcommand{\lrip}[1]{\mathopen{}\left\langle#1\right\rangle}
\newcommand{\Lrip}[1]{\mathopen{}\big\langle#1\big\rangle}
\newcommand{\LRip}[1]{\mathopen{}\Big\langle#1\Big\rangle}

\renewcommand{\t}[1]{\smash{\tilde{#1}}}
\newcommand{\wt}[1]{\smash{\widetilde{#1}}}
\newcommand{\wh}[1]{\smash{\widehat{#1}}}
\renewcommand{\O}{O}
%\renewcommand{\O}{\mathcal{O}}
\newcommand{\OO}[1]{\O\lr{#1}}
\newcommand{\tO}{\wt{\O}}
\newcommand{\tTheta}{\wt{\Theta}}
\newcommand{\E}{\mathbb{E}}
\newcommand{\EE}[1]{\E\lrbra{#1}}
\newcommand{\var}{\mathrm{Var}}
\newcommand{\tsum}{\sum\nolimits}
%\newcommand{\trace}{\mathrm{tr}} % declared as mathop
\newcommand{\ind}[1]{\mathbb{I}\lrset{#1}}
\newcommand{\non}{\nonumber}
\newcommand{\poly}{\mathrm{poly}}

% \newcommand{\tr}{^{\mathsmaller\top}}
\newcommand{\tr}{^{\mkern-1.5mu\scriptstyle\mathsf{T}}}
\newcommand{\pinv}{^{\dagger}}
\newcommand{\st}{\star}
\renewcommand{\dag}{\dagger}

%\newcommand{\kl}[2]{\mathsf{D}_\mathrm{KL}(#1,#2)}
\newcommand{\tv}[2]{\mathsf{D}_{\mathrm{TV}}(#1,#2)}
\newcommand{\ent}{\mathsf{H}}
\newcommand{\info}{\mathsf{I}}

\newcommand{\reals}{\mathbb{R}}
%\newcommand{\eps}{\varepsilon}
\newcommand{\eps}{\epsilon}
\newcommand{\sig}{\sigma}
\newcommand{\del}{\delta}
\newcommand{\Del}{\Delta}
\newcommand{\lam}{\lambda}
\newcommand{\half}{\frac{1}{2}}
\newcommand{\thalf}{\tfrac{1}{2}}

\newcommand{\eqdef}{\stackrel{\text{def}}{=}}
\newcommand{\eq}{~=~}
\renewcommand{\leq}{~\le~}
\renewcommand{\geq}{~\ge~}
\newcommand{\lt}{~<~}
\newcommand{\gt}{~>~}
\newcommand{\then}{\quad\Rightarrow\quad}

\let\oldtfrac\tfrac
\renewcommand{\tfrac}[2]{\smash{\oldtfrac{#1}{#2}}}


\let\nablaold\nabla
\renewcommand{\nabla}{\nablaold\mkern-1mu}

%\let\cdotold\cdot
%\renewcommand{\cdot}{\mkern-2.5mu\cdotold\mkern-2.5mu}
%\newcommand{\param}{\cdotold}

